\subsection{与正型函数的关系}

\begin{enumerate}
\def\labelenumi{(\arabic{enumi})}
\setcounter{enumi}{65}
\tightlist
\item
\end{enumerate}

命题 11. --- 设 E 是有限维向量空间，\(\mu\) 是 E 上的有界（正）测度，\(\Phi\) 是 \(\mu\) 的傅里叶变换。函数 \(\Phi\) 在 \(E'\) 上连续且为正型函数。

由第~8号的命题 9 可得 \(\Phi\) 的连续性。

证明 \(\Phi\) 是正型的。令 \(x_1', \ldots, x_p'\) 属于 E'，\(c_1, \ldots, c_p\) 为复数。则

\[
\begin{array}{r l} \sum_ {j, k} c _ {j} \overline {{{c}}} _ {k} \Phi (x _ {j} ^ {\prime} - x _ {k} ^ {\prime}) & = \int_ {\mathrm{E}} \sum_ {j, k} c _ {j} \overline {{{c}}} _ {k} e ^ {i \langle x, x _ {j} ^ {\prime} - x _ {k} ^ {\prime} \rangle} d \mu (x) \\ & = \int_ {\mathrm{E}} \left| \sum_ {j = 1} ^ {p} c _ {j} e ^ {i \langle x, x _ {j} ^ {\prime} \rangle} \right| ^ {2} d \mu (x) \geqslant 0. \end{array}\tag{Q.E.D.}
\]

可以证明一个称为博赫纳定理的逆命题：\(E'\) 上每个连续的正型函数都是某个有界（正）测度的傅里叶变换。\(^{*}\) 在第~12号余下部分假设这一结果成立。

定理 4. --- 设 E 是局部凸空间。傅里叶变换是从 E 上的投影测度集合到  上正型函数集合的双射，其中后者的函数限制在每个有限维子空间上都连续。

我们知道（第~3号的命题 3），傅里叶变换是单射。令 \(\mu = (\mu_{\mathrm{V}})_{\mathrm{V}\in \mathcal{F}(\mathbf{E})}\) 是 \(\mathbf{E}\) 上的投影测度，\(\Phi\) 是其傅里叶变换。令 \(\mathbf{T}\) 是 \(\mathbf{E}'\) 的有限维子空间，\(\mathbf{V}\) 是 \(\mathbf{T}\) 中 \(\mathbf{E}\) 的正交补。可将 \(\mathbf{T}\) 认作 \(\mathbf{E} / \mathbf{V}\) 的对偶；\(\Phi_{\mathrm{T}}\) 是 \(\Phi\) 在 \(\mathbf{T}\) 上的限制，也是 \(\mu_{\mathrm{V}}\) 上有界测度 \(\mathbf{E} / \mathbf{V}\) 的傅里叶变换。由命题 11，\(\Phi_{\mathrm{T}}\) 在 \(\mathbf{T}\) 上连续且为正型。由于 \(\mathbf{T}\) 任意，显然 \(\Phi\) 在 \(\mathbf{E}'\) 上为正型。

反过来，令 \(\Phi\) 是 \(\mathbf{E}'\) 上的正型函数，且其限制在 \(\mathbf{E}'\) 的每个有限维子空间上都连续。对于每个 \(V \in \mathcal{F}(\mathbf{E})\)，将 E/V 的对偶认作 \(E/V\) 中 V 的正交补 \(V^{\circ}\)；\(V\)\(\mathbf{E}'\)\(\Phi_V\) 是 \(\Phi\) 在 \(V^{\circ}\) 上的限制，它连续且为正型，因此由博赫纳定理，存在 E/V 上的有界（正）测度 \(\mu_V\)，其傅里叶变换为 \(E/V\)\(\Phi_V\)。设 \(V\) 和 \(W\) 属于 \(\mathcal{F}(\mathbf{E})\) 且 \(W \subset V\)，并令 \(p_{\mathrm{VW}}\) 为从 E/W 到 E/V 的典范映射；在上述认同下，\(E/W\)\(E/V\)\(^{t}p_{\mathrm{VW}}\) 是从 \(V^{\circ}\) 到 \(W^{\circ}\) 的嵌入。由第~3号的公式 (4)，于是有

\[
\mathcal {F} \big (p _ {\mathrm{VW}} (\mu_ {\mathrm{W}}) \big) = \big (\mathcal {F} \mu_ {\mathrm{W}} \big) \circ {} ^ {t} p _ {\mathrm{VW}} = \Phi_ {\mathrm{W}} \circ {} ^ {t} p _ {\mathrm{VW}} = \Phi_ {\mathrm{V}} = \mathcal {F} \mu_ {\mathrm{V}},
\]

因此由第~3号的命题 3，\(p_{\mathrm{VW}}(\mu_{\mathrm{W}}) = \mu_{\mathrm{V}}\)。所以族 \(\mu = (\mu_{\mathrm{V}})_{\mathrm{V} \in \mathcal{F}(\mathrm{E})}\) 是 E 上的投影测度；显然 \(\Phi\) 是 \(\mu\) 的傅里叶变换。

推论. --- 设 F 是桶状核空间；给 \(F'\) 赋予局部凸拓扑 T，使其介于弱拓扑 \(\sigma(\mathbf{F}',\mathbf{F})\) 与 F 的紧凸子集上的一致收敛拓扑之间。傅里叶变换是从 \(F'\) 上的有界（正）测度集合到 F 上连续正型函数集合的双射。

这立即由第 4 号定理和第~10号定理 2 的推论得出。*

